本研究針對語音深偽偵測（Audio Deepfake Detection）領域中日益嚴重的跨域泛化瓶頸，系統性地檢驗並比較當前主流的「預訓練—微調（Pre-train, Fine-tune）」範式（以 wav2vec 2.0 配合圖神經網絡後端的 W2V-AASIST，以及基於 Conformer 的 MFA-Conformer 為代表），以及以 Prompt Tuning 為核心的新興範式（以基於生成式語音語言模型 GSLM 之 SpeechPrompt v2 為代表）。透過橫跨 ASVspoof 2019/2021、InTheWild，以及多語言資料集 MLAAD 的大規模實驗，本研究深入剖析兩種範式在面對未知攻擊與真實環境下的防禦特性與先天限制。實驗結果揭示，雖然預訓練—微調範式憑藉豐富的底層連續聲學細節在同域測試中具備極高精準度，但在跨域場景中容易出現嚴重的效能衰退；相對地，Prompt Tuning 範式雖具備極優異的參數訓練效率（僅需調整約 0.13M 參數，遠少於 W2V-AASIST 的 317M），卻受前端聲學特徵離散化過程中不可逆的資訊流失所困，大幅限制了其偵測上限（同域 EER 由硬性量化之 18.50\% 經逐步鬆綁至 3.84\%，已大幅逼近預訓練—微調範式水準）。本研究透過設計軟性離散特徵與可學習群集中心等消融實驗，實證了離散化前處理瓶頸對 Prompt Tuning 範式之根本制約，並同時驗證了藉由圖神經網絡進行多層次特徵聚合，能有效紓解預訓練—微調模型的跨域效能飽和。本研究之實證結果為未來深偽偵測防禦系統的架構選擇與長遠演進方向，提供了關鍵且系統性的研發視角。此外，本研究開發之即時深偽語音偵測展示系統已公開於 GitHub，供後續研究參考與部署。

\begin{flushleft}
\mbox{{\bf 關鍵字}： 語音深偽偵測、跨域泛化、Prompt Tuning、語音語言模型}
\end{flushleft}

